% SOURCE docs/release/2026-10-29/source/docs/writing-repair/2026-10-29/src/bc7aff9d26aad5fb.tex CHAR 0
\chapter{Defining economic variables and measuring them}
\label{exec17:ontology}

The application examples leave two connected questions: what the records
measure, and what the available evidence permits us to assert about them.
This chapter develops both. It begins with constructs, measurement error,
units and clocks; compares ontology and annotation tools; then examines
conflicting expert evidence and joint graph advice. The final Bayesian study
shows why response frequencies and accepted graph probabilities must remain
distinct. Each step asks what new information an assisted workflow supplies.

A servicing proposal may ask whether an early support offer reduces financial
distress over the next thirty days. The bank can observe payment events,
balances, customer messages and employees' notes, but none of these is simply
the construct of financial distress. A payment can be made by drawing on
another account; a note can describe anticipated hardship rather than a
realized missed payment; an employee can write more carefully after making an
offer. An assistant that assigns all these records the same variable name
changes the observation on which the causal question depends.

The proposed ontology layer has a specific purpose: make the relationship
between a construct, its operational measure and its original observation
explicit enough to challenge. This can help a scientist formulate a question,
a data engineer construct a reproducible observation and a reviewer locate a
disputed assumption. It does not determine which construct the bank should
value, establish the truth of a causal relation, or make an unobserved
confounder observable. The distinction has consequences even when the
intervention itself is randomized.

\section{From a business construct to a defined estimand}

Let \(A_i\in\{0,1\}\) denote a support offer at decision time zero for eligible
account \(i\). Define \(Y_i(a)\) as the indicator of a specified missed-payment
event by day thirty under intervention \(a\). With a fixed eligible population
and account weighting, the target is
\[
 \tau_Y=E\{Y_i(1)-Y_i(0)\}.
\]
Consistency requires a sufficiently specified offer, comparator and event;
random assignment and appropriate handling of interference would justify a
simple experimental comparison. An observational proposal additionally needs
a defensible adjustment or other identification argument. Calling a field a
confounder does not supply that argument.

Now let \(T_i(a)\) be a service note and \(L_i(a)=g_\theta(T_i(a))\) its binary
label under a frozen extraction rule \(g_\theta\). Even if the label is named
after missed payments, the experiment on labels targets
\[
 \tau_L=E\{g_\theta(T_i(1))-g_\theta(T_i(0))\}.
\]
This is a defined effect of the offer on a coded textual outcome. Equating it
with \(\tau_Y\) requires an additional measurement argument. If the desired
target is latent hardship \(U_i(a)\), neither an observed payment nor a
plausible description of it automatically identifies that target.

Egami and colleagues develop this distinction through a learned text mapping:
the representation helps define the causal outcome, and learning it from the
same treatment-assigned evaluation documents can induce dependence and
overfitting. Their train/test separation holds the learned mapping fixed for
effect estimation on the evaluation sample.~\citep{exec13_census_texts}
That remedy addresses reuse of evaluation information; it does not establish
that a learned category measures hardship without error. The earlier
measurement discussion's residual-confounding example makes a different
point: even a useful proxy can leave confounding and reverse a coefficient's
sign. Both problems survive a syntactically valid ontology.

Keith, Jensen and O'Connor explain why the causal role of text also matters:
text may measure confounding, mediate a treatment or be selected by a collider
mechanism.~\citep{exec13_census_text_review_arxiv} A note written after the offer
cannot be admitted as a baseline covariate merely because it mentions a
pre-existing condition. Its information might have been elicited by the offer.
The proposal must distinguish the underlying condition's event time, the
record's availability and the process that produced the record.

For each variable, the proposed representation retains five linked objects:
the business construct, the operational definition, the observation or
transformation, the role in a particular causal model, and the target
population's unit and weighting. These links can have different epistemic
statuses. A schema can verify that a payment record has a date and amount;
a reviewer can accept a definition for a study; a causal analyst can still
leave the identification assumptions unresolved. A single confidence score
cannot substitute for these separate decisions.

\section{What formal ontologies and language models can contribute}

Pressat-Laffouilhere's OntoBioStat thesis treats variable selection as a problem
of formal representation and supplied relations. Its account of recoding is
particularly relevant: changing a variable's categories creates a different
variable, whose context, possible values, missingness and transformation must
be represented. It also distinguishes hypothetical causal relations from
relations under validation and illustrates sensitivity analysis by retaining
or removing uncertain links before rerunning a reasoner.~\citep{exec17_ontology_thesis}
For a bank, an analogous distinction would separate a note mentioning
hardship, a customer reporting hardship and reviewed evidence supporting an
effect of hardship on payment. An inference explanation can show which
supplied relations were used. It cannot establish that those relations are
true, or guarantee a valid adjustment set for every causal graph.

The two knowledge-graph surveys offer complementary architectural choices.
Jeff Z. Pan and colleagues contrast explicit-knowledge-first with
parametric-knowledge-first approaches. The former can expose a curated
vocabulary and links for the assistant to retrieve; the latter can propose
entities and relations that subsequently need canonicalization and
review.~\citep{exec17_ontology_kg_challenges} Shirui Pan and colleagues organize
the broader interaction into knowledge-graph-enhanced language models,
language-model-augmented graphs and synergized systems. Their retrieval,
representation and graph/text fusion mechanisms explain how an assistant might
find a definition or propose a relation.~\citep{exec17_ontology_kg_roadmap}
The question-answering paths and attention scores in these architectures do
not identify the effect of an intervention.

A useful design keeps relation types separate. A type-of relation
belongs to a taxonomy; derivation describes data lineage; association
records an observed relationship; a hypothesized cause belongs to a
candidate causal model. A household owns accounts, but the ownership edge is
not a causal arrow from a household variable to an account outcome. A synonym
link can support a search without authorizing the merger of two measurements.
An ontology reasoner checks implications of its formal premises; a causal
identification procedure has different premises and a different output.

Tan and colleagues' annotation survey separates generation, assessment and
utilization of annotations, and discusses hallucination, feedback from
synthetic data and manual or automatic validation.~\citep{exec17_ontology_annotation}
This separation suggests a staged bank workflow: an extractor proposes a
field and cites the passage; a typed validator checks identity, unit and
clock; a reviewer adjudicates meaning; a causal analysis uses the resulting
measurement under explicit assumptions. Agreement among model-generated
labels is evidence about agreement, not independent evidence about the
underlying construct. Generated respondents are especially unsuitable as
substitutes for observed customers in an effect estimate.

Schwitter, Bach and Klamm distinguish text as outcome, treatment and control
and illustrate substantive measurement decisions in sociological
applications. In one reported annotation example, zero-shot kappa is
\(0.63\pm0.01\), compared with \(0.78\pm0.01\) under three-shot prompting,
across three reruns. These are agreement results in the authors' setting,
not calibrated bank measurement errors or causal confidence intervals.
Their distinction between annotating real respondents and generating
respondent simulations is directly relevant here. Their illustrated noise
experiment discusses a tendency toward type-II error; differential
measurement need not behave that way.~\citep{exec17_ontology_sociology}

\section{High label accuracy can coexist with a false monetary conclusion}

Consider an exact authored population with randomized, equally sized offer
and control arms. The true adverse-event risk is \(p_0=p_1=1/5\). Write
\[
 s_a=P(L=1\mid Y=1,A=a),\qquad
 f_a=P(L=1\mid Y=0,A=a).
\]
The observed label risk is
\[
 q_a=P(L=1\mid A=a)=s_ap_a+f_a(1-p_a)
     =(s_a-f_a)p_a+f_a.
\]
This identity follows by partitioning on the true outcome. It does not
assume that an observed label is the true outcome. Suppose the control arm has
\((s_0,f_0)=(4/5,1/20)\), while the offered arm has
\((s_1,f_1)=(9/10,3/20)\). For instance, the intervention may change employees'
documentation and the terms customers use. The stipulated numbers imply
\[
 q_0=\frac15,\qquad q_1=\frac3{10},\qquad q_1-q_0=\frac1{10},
 \qquad p_1-p_0=0.
\]
Control accuracy is \(0.92\), offered-arm accuracy is \(0.86\), and pooled
accuracy is \(0.89\). Despite that apparently respectable aggregate accuracy,
the label comparison reports a ten-percentage-point adverse effect under a
true null. At a stipulated loss of USD~1,000 per true event, treating labels
as events invents USD~100 of additional loss per account. These are
illustrative exact values, not estimates from customers or model
calls.~\citep{exec17_ontology_oracle}

If valid arm-specific measurement rates were known and \(s_a\ne f_a\), the
algebraic correction
\[
 p_a=\frac{q_a-f_a}{s_a-f_a}
\]
would recover \(1/5\) in each arm in this example. The rates must describe the
actual measurement process and target population. A representative validation
sample would introduce estimation uncertainty; small \(s_a-f_a\) amplifies
that uncertainty; shifts in language or workflow can invalidate the
calibration. Prompt confidence is not a supplied value of \(s_a\) or \(f_a\).

Without measurement information the problem is one of identification, not
just classifier improvement. The same observed pair
\((q_0,q_1)=(1/5,3/10)\) also arises under perfect measurement and true risks
\((p_0,p_1)=(1/5,3/10)\). Thus the observed treatment-label law is compatible
with true effects zero and \(1/10\). More copies of that law cannot choose
between these worlds. A useful proposal can compare sensitivity ranges,
identify an additional validation measurement, or explicitly target the
effect on coded notes. It cannot recover a unique latent effect by renaming
the label.

This example explains why validation must be stratified by scientifically
relevant circumstances: treatment arm, language, document source and time
can alter error rates. The eventual sample design must justify how
adjudicated observations represent the target population. Such expensive
validation follows a promising prototype; the proposal can establish the
failure mechanism and candidate remedies using generated laws now.

\section{Entity, unit and time are part of the variable}

Aggregation can change the question while preserving every individual number.
Suppose two households have stipulated household effects of USD~\(10\) and
USD~\(-10\), but nine and one accounts respectively. The equal-household mean
is zero. Replicating each household's value once per account gives
\[
 \frac{9(10)+1(-10)}{9+1}=8.
\]
This is an account-weighted \emph{household} effect, not proof of an
account-level causal mechanism. Ten rows still refer to two households.
The representation must retain entity identifiers, grouping links,
the level of treatment assignment and the target weighting. A precision
calculation cannot treat repeated household records as independent households.
\citep{exec17_ontology_oracle}

Units require more than a monetary label. Monthly debt service of
USD~1,500 and monthly income of USD~3,000 imply a debt-service ratio of \(1/2\).
Dividing the same monthly obligation by annual income of USD~36,000 gives
\(1/24\). Converting annual income to monthly income restores \(1/2\).
The conversion is justified by an explicit common-period definition in this
example. It would not reconcile gross with net income, household with
individual income, a stock with a flow, or average with end-period balances.
Those are different constructs or aggregation rules. Currency conversion
would additionally require the rate and valuation date.

An observation's event time and usable time also differ. A fact about day
one that becomes available on day five cannot inform a day-three decision.
For a derived value computed from inputs with availability times two and
five, availability cannot precede five. A declared time of four fails even
when the input events occurred earlier. Actual production delay may make
availability later still. These examples define a historical information
constraint; they do not prohibit outcomes from becoming observable after
the treatment decision.~\citep{exec17_ontology_oracle}

The local proposal implementation already checks declared measured/latent
types, derived lineage, event-before-availability order, input availability
and information at a decision. Its restricted static-effect projection
retains variable identity, role and unit, keeps full declarations as evidence
and refuses a mismatch between outcome and loss units because it has no
utility transformation. These are useful declaration checks. They neither
audit individual rows nor supply a verified extraction, household linkage or
currency-conversion service.~\citep{exec17_ontology_sources}
The proposed next layer must carry record-specific clocks and provenance
through the conversion, and preserve an explicit unresolved state when
availability or identity cannot be established.

\section{Missing notes and imputed values need an observation model}

Even perfect annotation of the notes that exist can mislead. Let
\(R_i=1\) indicate that a note is observed. In another exact authored law, the
true adverse-event risk is \(1/5\) in both randomized arms. All control notes
are observed. In the offered arm, all adverse events have a note, but only
one quarter of non-events do. Then
\[
 P(R=1\mid A=1)=\frac15+\frac45\frac14=\frac25,
 \qquad
 P(Y=1\mid A=1,R=1)=\frac{1/5}{2/5}=\frac12.
\]
Comparing observed notes gives an apparent contrast \(1/2-1/5=3/10\)
under a true null. Selection depends jointly on treatment and outcome.
Annotation accuracy among observed notes cannot repair the missing
population.~\citep{exec17_ontology_oracle}

Known positive observation probabilities would permit the illustrated
weighted population calculation to recover \(1/5\); those probabilities are
stipulated in this reference, not estimated from notes with unobserved
outcomes. If the relevant selection law is unknown, collecting representative
outcomes or adding justified assumptions is necessary. An assistant can expose
the missingness mechanism, propose a sensitivity analysis and keep a value
marked missing. Filling it with a plausible sentence does not reveal what
the customer would have reported.

This caution is consistent with the primary quantitative-knowledge study
reviewed earlier. Its imputation comparisons report that KNN and random
forest methods outperform the examined LLM alternatives in the stated
settings; ambiguous descriptions and units are discussed as contributors to
error. Recognition of familiar datasets also raises contamination concerns.
\citep{v12_la_2402_07770} These findings support retaining ordinary imputation
and dictionary baselines. They do not imply that all LLM extraction fails,
or that improved prediction identifies the causal effect when observation
depends on an unobserved outcome.

The ontology should represent missing, not applicable, not collected,
redacted, delayed and model-imputed values distinctly. It should retain the
original observation indicator even if an imputation is created. A changed
imputation model creates a new derived measurement revision; it does not
rewrite the original record or convert a latent quantity into an observed
fact. Inference about imputed values needs a separate justification of the
imputation and sampling laws.

\section{Benchmark evidence and a concrete package decision}

OntoURL makes the distinction between understanding, reasoning and generation
testable through fifteen ontology tasks. Its construction uses rule
extraction, a reasoner for reasoning tasks, distractors and quality review.
The inspected preprint reports strong performance on some recognition tasks
and weak performance on several structured generation tasks. For example,
its Qwen2.5-72B row gives class-definition and class-relation understanding
(U1 and U2) accuracies of \(89.1\%\) and \(92.6\%\), and rule-based and
description-logic reasoning (R4 and R5) accuracies of \(75.6\%\) and \(68.4\%\).
Class-hierarchy, property-relation and constraint construction (L2--L4) have
Triple-F1 scores of \(0.1\), \(0.0\) and \(1.0\); ontology alignment (L5)
has Tuple-F1 of \(21.6\), on the table's percentage scale.
These are different tasks and metrics, not a causal-graph accuracy
scale.~\citep{exec17_ontology_ontourl}

The same version contains a material reporting inconsistency: Table~2 gives
L1 BERTScores around 79--83, while nearby prose says they are typically below
10 and Table~3 uses another apparent scale. No aggregate comparison should
silently reconcile these numbers. The three PhD-level human participants
answer thirty sampled questions per task, while the table's model results
cover the full dataset. This also limits a simple human-versus-model
ranking. Prompting effects vary across tasks, and the authors acknowledge
restricted domains, incomplete task coverage and an English-only scope.
The usable lesson is to evaluate the actual task and output semantics;
the disputed scale and sampling differences require separate resolution.

OntoLearner offers a more concrete reuse opportunity: a Python framework
with task-specific data and evaluation paths. Its paper reports
domain-dependent error patterns and calibration behavior, and itself argues
that F1 alone misses important failure modes.~\citep{v12_la_2607_01977}
The inspected revision's document loader retains document identifiers and
extraction methods, its evaluator dispatches among four task metrics, and
its non-taxonomic metric only treats three declared relation names as
symmetric. In particular, the code does not symmetrize every relation:
reversing a causal triple is different from reversing a
\texttt{sameas} triple.~\citep{exec17_ontology_sources}

There is a specific reason to adapt its evaluation interface before any
causal use. Its text-to-ontology metric concatenates the document identifier
and extracted value, then greedily matches token sets using a Jaccard
threshold of \(0.8\). Take two different document identifiers and the
same nine-word phrase, \emph{monthly household income before tax during the
reporting period}. The token intersection has size nine and the union
eleven, so similarity is \(9/11>0.8\) even though the document link is wrong.
This is an exact authored counterexample to the inspected scoring
definition, not an executed upstream-package result.
\citep{exec17_ontology_sources,exec17_ontology_oracle}

A causal observation contract should first require exact document identity,
then compare meaning within that document; source-span identity, units and
clocks need their own checks. Lexical similarity can remain a useful
diagnostic after those requirements. Similarly, a taxonomy metric can measure
hierarchical agreement without validating an intervention graph, and
deduplication can alter denominators if the intended unit is an observation
rather than a unique term.

The package decision is to use OntoLearner as a reference and candidate
benchmark component, with an adapter proposed for typed observations and
separate semantic scoring. Direct adoption of the inspected metrics as a
causal acceptance gate is rejected. The inspected downloads do not pin
dataset revisions; enabled ontology-import failures can be logged while
processing continues; missing dataset splits can also continue with warnings.
A reproducible evaluation must freeze required dependencies and splits and
explicitly fail or mark incomplete required missing inputs. The base
learner's formatting path does not establish all concrete subclass behavior.
Its MIT license permits relevant code reuse subject to the notice, but does
not establish rights to each external ontology, dataset, model or dependency.
Full runtime qualification remains a separate proposed step.
\citep{exec17_ontology_sources}

\section{A bounded public and generated comparison before bank validation}

The proposed prototype takes a document or record, a vocabulary revision, a
business question and its decision clock. It returns candidate observations
with the original record and span, entity and aggregation level, construct and
operational definition, value domain and unit, event and availability clocks,
transformation lineage, extraction revision, and an explicit evidence status.
Candidate causal roles and relations are separate objects linked to the
model in which they are proposed. Unresolved meaning remains unresolved.
An assistant's numeric confidence may be stored as a model output, but is
not interpreted as a calibrated measurement-error probability.
\citep{exec17_ontology_contract}

The initial comparison should contain a deterministic dictionary/schema
baseline, an ontology-assisted mapping alternative, and a constrained
language-model extractor if a later permitted model trial is warranted.
The six exact families above supply tractable semantic reference cases.
Additional generated variants should include ambiguous entities, negation,
incompatible definitions, conflicting passages, renamed identifiers and
unavailable facts. Positive cases establish what a valid conversion must
retain; negative cases establish the appropriate refusal or unresolved
output. A case with observationally equivalent causal worlds should reward
preserved ambiguity, not a lucky choice of the generator's hidden world.

Public ontology benchmarks can test recognition and formal relation tasks;
their distributions, rights and score definitions must be inspected before
use. Public treatment text and code identified by the sociology article are
possible comparison materials, whereas the article's onsite-only respondent
text is not a public dataset. None supplies bank causal truth. Generated
banking records can isolate known unit, clock, grouping and selection
mechanisms without pretending to be authentic bank observations. They also
allow newly constructed combinations that reduce straightforward lookup,
although they do not establish absence of contamination or real-world
generalization.

Before model comparison, freeze the case-generating rules, development and
evaluation split, document families, source revisions, permitted tools,
prompts, model identity, budgets and stopping rule. Keep paraphrases and
records from the same entity together when defining evaluation units. Freeze
the measurement mapping before effect evaluation. Report exact identity,
unit and clock preservation; typed relation direction; correctness of
acceptance, refusal and unresolved dispositions; performance on all requests
and on accepted outputs; and effects on the downstream target. Retain
malformed responses and missing documents in the denominator. Human
adjudication of meaning, if introduced later, needs an independent label
procedure and disagreements must remain visible.

Passing all stipulated semantic cases is a bounded engineering criterion,
not a guarantee of future causal validity. The comparison should separately
measure whether assistance saves review effort after retrieval, correction,
adjudication, retries and maintenance. A proposed economic calculation can
compare avoided review and error loss with these additional costs; it cannot
assign actual bank savings before the relevant error frequencies, prices
and consequences are measured. The next investment decision is whether a
permitted prototype demonstrates enough benefit to justify expensive
representative measurement validation.

The remaining research is specific. Historical availability and entity
matching need evidence for actual records; latent constructs need measurement
assumptions and representative validation; selection and imputation need a
defensible observation law; causal sufficiency needs additional scientific
assumptions or information. Better extraction can reduce some implementation
errors, but observationally equivalent measurement mechanisms require new
information or an explicitly broader answer. The proposed program can develop
and compare these remedies now without claiming that an ontology, benchmark
score or language model has settled them. The following studies ask how to
handle the evidence used to define those measurements and restrict their
causal models: how to distinguish textual knowledge from observations, retain
expert disagreement, and use advice without treating repeated agreement as
independent confirmation.
\label{exec17:ontology-end}


% SOURCE docs/release/2026-10-29/source/docs/writing-repair/2026-10-29/src/28e65a00a5e889d5.tex CHAR 0
\section{From language-model knowledge to valid causal observations}
\label{sec:exec24-css-knowledge}
A bank asking a language model to interpret a complaint is performing measurement before it performs causal inference. The model might classify financial distress, extract a promised payment date or propose an explanation for a servicing outcome. Each output has a different relation to the scientific question. A useful proposal must specify the measured construct, its time of availability, its error mechanism and the permitted role of the model in subsequent identification. The computational-social-science study by Ziems and colleagues and the knowledge-mechanism survey by Wang and colleagues clarify these boundaries. They support a disciplined assistance workflow while leaving causal validity and operational qualification to separate evidence.~\citep{exec24_css_original,exec24_knowledge_original,exec24_css_knowledge_reference}

\subsection{What the computational-social-science evaluation establishes}
The inspected Ziems text is the genuine 54-page author-uploaded version dated February 2024, carrying Computational Linguistics volume and publication information. Its evaluation covers 13 historical language models and 25 English benchmarks, including utterance, conversation and document classification and generation tasks. Test construction uses class-stratified samples of at most 500 items per task; without an original designated test split, the sample comes from the available dataset. Classification uses macro F1 averaged over five prompt variants. Supervised baselines receive task-specific tuning and multiple random seeds. This is a practical utility comparison rather than a controlled isolation of one modeling ingredient.~\citep{exec24_css_original}

The paper reports stronger fine-tuned baselines overall for classification, alongside useful agreement with human labels for some zero-shot tasks. Its error analysis is consequential: mistakes can coincide across models, performance varies across classes, and few-shot examples do not improve every task consistently. Generation quality uses different human scoring and ranking designs. Some assessments use recruited domain experts; the ranking procedure uses blinded author judgments and unanimity. These procedures measure declared judgments, not interchangeable forms of truth. Favorable explanation ratings do not establish a correct causal explanation.~\citep{exec24_css_original}

The tasks are not exhaustive, cultural and linguistic coverage is bounded, systematic demographic performance analysis was not feasible, and training-data overlap remains unresolved. The paper explicitly separates predictive characterization of text from causal explanations. Its historical model sizes, architectural descriptions and prices are not current specifications or independently verified facts in this proposal. We rely on inspected experimental designs, reported task-specific findings and limitations, without reproducing its model runs or claiming contemporary superiority.~\citep{exec24_css_original}

\subsection{Translate the business question into a measurement contract}
Consider whether a revised servicing message reduces subsequent financial distress. Before proposing a diagram, define the unit, intervention, horizon, distress construct, observation source and monetary decision. A complaint-based label is one observation of distress. The message may change whether a customer complains, how the complaint is phrased or how staff record it. The diagram should distinguish the underlying outcome $Y$, treatment $A$, observed text $T$, generated label $\widetilde Y$ and the indicator for availability of text. Omitting these nodes can turn a measurement procedure into an unsupported causal assumption.~\citep{exec24_css_knowledge_reference}

A model can propose operational definitions, alternative readings, timestamps and passages supporting candidate variables. The scientific owner then determines whether those readings preserve the estimand and whether an edge represents a temporal fact, substantive assumption, statistical hypothesis or unresolved possibility. Retain graph alternatives when evidence cannot resolve them. A fluent explanation or majority vote over paraphrases is insufficient to promote an edge to an identified causal relation.~\citep{exec24_css_knowledge_reference}

An authored binary example shows why aggregate accuracy is inadequate. Let true distress prevalence be $p_1=p_0=0.20$. Define arm-specific sensitivity $s_a=P(\widetilde Y=1\mid Y=1,A=a)$ and specificity $c_a=P(\widetilde Y=0\mid Y=0,A=a)$. Observed label prevalence is
\[
 q_a=s_a p_a+(1-c_a)(1-p_a).
\]
With $(s_1,c_1)=(0.90,0.90)$ and $(s_0,c_0)=(0.70,0.95)$, both arms have 90 percent accuracy, but $q_1=0.26$ and $q_0=0.18$. The labels imply an eight-point contrast when the true prevalence contrast is zero. A stipulated 1,000 USD loss per distress event converts this to an 80 USD distortion per customer in the corresponding loss contrast. This is a possible measurement failure, not an observed model or bank result.~\citep{exec24_css_knowledge_reference}

If target-population sensitivity and specificity are known and $s_a+c_a\ne1$, inversion gives
\[
 p_a=\frac{q_a+c_a-1}{s_a+c_a-1}.
\]
This is a conditional route. An audit must identify the error rates, preserve label-inclusion weights and uncertainty, and address shifts across treatment, language, time and customer groups. Near a zero denominator, prevalence is weakly informed; at zero, this equation does not identify it. Interpreting the repaired prevalence contrast as a treatment effect still requires assignment and observation assumptions.~\citep{exec24_css_knowledge_reference}

\subsection{Prompt variation, refusal and dependent evaluation}
The Ziems study averages prompt perturbations to examine performance variation. In the proposed workflow, repeated prompts remain related measurements unless a justified law establishes independence. If five scalar errors have variance one and pairwise correlation $\rho$, then
\[
 \operatorname{Var}(\overline E)=\frac{1+4\rho}{5}.
\]
At stipulated correlation 0.80, the variance is 0.84 rather than 0.20 under independence. The equivalent independent count for this variance comparison is $5/(1+4\rho)=25/21$, about 1.19. This exact covariance illustration is not an estimated correlation for the paper's models.~\citep{exec24_css_original,exec24_css_knowledge_reference}

Refusal changes the experiment. Historical constrained-output prompts facilitate fixed-taxonomy scoring, but a causal product should abstain on underidentified, contradictory or unauthorized questions. Predeclare when refusal is correct and report answer quality and availability separately. In an authored fixture with 100 requests, 90 answers and 81 correct answers, answered-only accuracy and coverage are each 90 percent, while correct answers per request are 81 percent. Refusals need a reason and correctness disposition; neither dropping them nor automatically treating them as successful analysis is defensible.~\citep{exec24_css_original,exec24_css_knowledge_reference}

Hold out identities and time periods when repeated text or policy regimes could leak across development and evaluation. Separate prompt/model selection, taxonomy revision and final assessment. Label disagreement may indicate model error, construct ambiguity or deficient reference labels. Human adjudication has a cost and an inclusion mechanism; selectively checking only conspicuous mistakes does not identify a population error rate.~\citep{exec24_css_knowledge_reference}

\subsection{Knowledge mechanisms and the target of an intervention}
The inspected Wang survey is a genuine 39-page December 2024 author version. Published-work identity is corroborated by title and author sequence, without asserting publisher-byte equivalence. It organizes knowledge utilization into memorization, comprehension/application and creation, and considers individual and group evolution. Observation methods include probes, logit-based inspection and sparse representations; intervention methods include clean, corrupted and restored runs to locate components involved in an output. These study the computational system, not automatically the external world described by its text.~\citep{exec24_knowledge_original}

Let $F_\theta(x)$ be a model's answer to prompt $x$, with internal activation $H$. Intervening on $H$ asks how a computational change affects the answer. A bank's target instead concerns outcome $Y$ under business action $A$. Connecting the two requires evidence about representation, measurement and the external generating process. Activation restoration alone supplies no such bridge.~\citep{exec24_css_knowledge_reference}

For a finite counterexample, let $U$ be Bernoulli with probability one half and $X=U$. In world one, $Y=X$; in world two, $Y=U$. Both yield observed pairs $(0,0)$ and $(1,1)$ with equal probability. Yet the average effect of setting $X$ from zero to one is one in the first world and zero in the second. A model given the shared observations can produce identical answers and internal patterns in both worlds. Stipulated clean, corrupted and restored answer probabilities $0.90,0.20,0.80$ give restoration score $(0.80-0.20)/(0.90-0.20)=6/7$ in either world. Restoring an answer has not selected the external causal law.~\citep{exec24_css_knowledge_reference}

The survey cautions that mechanistic findings depend on method, metric and implementation. Its accounts of particular tools and datasets are discovery leads here; their original technical texts and source code require separate inspection before implementation claims. Mechanistic analysis may diagnose an explanation pipeline or evidence sensitivity while graph-equivalence, omitted-variable and intervention-design obligations remain.~\citep{exec24_knowledge_original,exec24_css_knowledge_reference}

\subsection{Retrieval, memory conflict and agent agreement}
The survey distinguishes conflict inside parametric memory from conflict between supplied context and stored knowledge, and discusses agent knowledge sharing and conformity. These categories help design tests. Its explanations of fragility and future inaccessible knowledge remain hypotheses, not general theorems. Limitations include unresolved hypotheses, omitted time/event domains, incomplete technical coverage and uncertain generalization across model scales. A survey taxonomy must not become a product guarantee.~\citep{exec24_knowledge_original}

Each retrieved banking statement should carry source identity, effective date, institutional scope and supported claim. A new policy document may not apply to a historical causal study. The application must resolve scope and time before accepting a proposed supersession. Conflicting sources should yield a documented alternative, acquisition request or refusal. Several agents repeating one source do not provide independent corroboration, and editing model memory does not establish truth.~\citep{exec24_css_knowledge_reference}

A public experiment can test correct evidence with misleading prompts, stale evidence with plausible answers, conflicting documents, missing dates, unsupported edges, duplicate sources and deliberately unidentifiable graphs. These cases have explicit semantic oracles before a bank purchases costly domain labels. Success establishes the tested contract, with transport to bank populations still requiring justification.~\citep{exec24_css_knowledge_reference}

\subsection{Executable proposal boundaries and complete cost}
The proposed CLI/MCP workflow separates question interpretation, retrieval, measurement, diagram proposals, identification and estimation. Each stage consumes immutable named artifacts and returns assumptions, alternatives, provenance and unresolved questions. A diagram needs definitions and clocks; identification needs the graph class and allowed assumptions; measurement needs the calibration population and abstention rule. The estimator should refuse a missing identification contract rather than infer one from prose. This is a proposed interface, not a deployed multi-repository integration.~\citep{exec24_css_knowledge_reference}

Known next steps are to select versioned models and public tasks, inspect the methods/packages actually used, freeze semantics, measure selection/dependence/refusal and compare complete held-out workflows. Mechanistic analysis can supplement these diagnostics when supported. Hidden assignment, unrestricted confounding, unjustified transport and absent institutional truth remain information gaps that more prompts alone cannot close. Authentic owners and later bank evidence follow a useful prototype.~\citep{exec24_css_knowledge_reference}

For the 100-request fixture, stipulate 240 USD total cost across model use, retrieval, adjudication and engineering. Cost is 2.40 USD per request, $8/3$ USD per answer and $80/27$ USD per correct artifact. These are explanatory denominators, not prices or measured returns. Compare them with decision loss, uncertainty and correct refusal while distinguishing historical findings, authored derivations and future implementation evidence.~\citep{exec24_css_knowledge_reference}


% SOURCE docs/release/2026-10-29/source/docs/writing-repair/2026-10-29/src/a133b8274089ca83.tex CHAR 0
\hypertarget{expert-evidence-disagreement-and-causal-refusal}{%
\section{Expert evidence, disagreement and causal
refusal}\label{expert-evidence-disagreement-and-causal-refusal}}

\label{exec14:expert-evidence}

A collections manager asks whether a new contact offer will increase
repayment within thirty days. An experienced officer says that offers
are allocated to accounts thought most likely to respond. Another
describes assignment as effectively random within the day's queue. A
retrieved operational note describes the standard queue but predates a
change in the assignment system. An assistant can locate those accounts,
identify the disagreement and make its consequences explicit. Counting
two statements in favor of one story does not resolve which assignment
mechanism operated for the accounts under study. This distinction
determines what the proposed product should accept from experts,
documents and language models.

The useful output is a collection of inspectable claims with their
supporting and opposing evidence, the assumptions still being made, and
the decisions that depend on them. A causal diagram becomes one
representation of those claims. The central design problem is how to
retain a plausible alternative when a passage is missing, an expert
account conflicts with a log, or two documents repeat the same original
assertion. The literature supplies useful components for this task, with
different evidence boundaries.

\hypertarget{what-retrieval-elicitation-and-interactive-auditing-establish}{%
\subsection{What retrieval, elicitation and interactive auditing
establish}\label{what-retrieval-elicitation-and-interactive-auditing-establish}}

Zhang et al.'s LACR constructs a skeleton using model background
knowledge, retrieved documents and optionally statistical discovery
output, before orienting the retained edges. Its experiments support a
conditional role for retrieval: the ASIA background-only skeleton has
reported F1 of 1, while adding documents lowers it to approximately
0.727; other benchmark cells improve. The orientation score conditions
on true-positive skeleton edges and must not be presented as whole-graph
accuracy. More retrieved text is therefore an input whose relevance and
reliability need evaluation. \citep{v12_la_2402_15301}

The algorithm also reveals an uncertainty problem that a superficial
reading of its prompts would miss. Individual responses may be UNKNOWN
and contribute zero to an aggregation score. Nevertheless, the final
rule deletes an edge when that score is nonpositive. If every response
is unknown, the edge is deleted. Retaining an unknown token in a prompt
is thus insufficient: the aggregation and graph-conversion stages must
preserve its meaning. The banking proposal should retain an unresolved
adjacency when evidence was not located, unless a separately justified
statistical or structural argument excludes it.
\citep{v12_la_2402_15301}

Soft elicitation offers a different route. Darvariu et al.~query both
variable orders, average the corresponding judgments and use the
resulting priors to bias graph search. Combined mutual-information and
LLM priors improve selected benchmark comparisons, but LLM-only priors
are not consistently beneficial; mutual information alone is strongest
on the specialized Child benchmark. This motivates sensitivity to the
influence assigned to elicited knowledge. It does not convert the
model's stated probability into a calibrated probability of a banking
edge. Selby et al.~provide complementary adverse evidence: some elicited
priors conflict with data, and ordinary nearest-neighbor and
random-forest imputations provide strong alternatives. Their discussion
of ambiguous units and recognizable public datasets is directly relevant
to financial data dictionaries and evaluation contamination.
\citep{v12_la_2405_13551,v12_la_2402_07770}

Interactive auditing can make suggestions easier to inspect. The early
\emph{Causal Auditor} organizes ratings, proposed confounders and
mediators into visual summaries. It reports 103 correct directions out
of 110 queries in its selected setting. That is a literature observation
from early cases, with no new evidence about collection offers. Its
reported BIC change also needs caution: the text calls a change from
-3673.4 to -6270.3 an improvement while the captions say higher scores
are better. Without a reconciled convention and reproduction, the
claimed percentage should not support the product's value case. In any
event, a change in model fit alone would not validate a causal
direction. \citep{exec14_expert_auditor}

The later CausalChat study reports incorrect-edge counts of 5/18 with
BIC feedback, 7/18 with LLM text, and 1/18 with its integrated
interface. Six students completed the stages in a fixed sequence and
selected different edges as they proceeded. These observations motivate
usability hypotheses but do not isolate the causal effect of assistance
under randomized task allocation. Expert feedback explicitly identifies
missing citations and uncertainty about the scientific evidence behind
the model's explanations. A persuasive display must therefore preserve a
route back to the primary passage and its limits.
\citep{v12_la_2410_14146}

Sng et al.~show why a metric also needs a precise interpretation. Their
hallucination criterion concerns contradictions across causal-rating
prompts. The retrieval experiment supplies a textual, mocked true graph,
and the paper explicitly cautions that reduced contradiction does not
establish accurate causal prediction. The proposed banking evaluation
should separately measure retrieval, statement support, graph
admissibility and intervention error; success on one cannot silently
substitute for the others. \citep{exec13_hallucination}

\hypertarget{formal-disagreement-handling-still-depends-on-its-premises}{%
\subsection{Formal disagreement handling still depends on its
premises}\label{formal-disagreement-handling-still-depends-on-its-premises}}

Li and Russo's 2026 ABAPC-LLM supplies a more explicit combination of
semantic and statistical constraints. It elicits required and forbidden
directions, parses them into a schema, and retains their intersection
across five model responses. Conditional-independence results enter a
causal argumentation framework; its solver seeks an admissible extension
after relaxing selected constraints. This makes the provenance and
consequences of a conflict more visible than an unrecorded majority
vote. \citep{exec14_expert_argumentation}

The implementation described by the paper is more restrictive than the
general idea of defeasible knowledge. A preliminary skeleton reduction
removes edges that are never reintroduced. Required LLM arrows are
discarded when their adjacencies have already been removed; the
remaining required and forbidden directions are enforced as hard
constraints. A stable output therefore establishes consistency relative
to the surviving restrictions. It cannot recover an edge excluded by an
erroneous preliminary decision, nor guarantee that a repeatedly
suggested hard constraint is true. The paper acknowledges this tradeoff
between tractability and reconsideration of the hard constraints.
\citep{exec14_expert_argumentation}

The study uses 54 generated graphs with five, ten or fifteen nodes and
repeated samples from their stipulated distributions. It reports
structural gains over several statistical and language-model
alternatives. The evaluation is useful because it exposes both data
constraints and semantic constraints, but its metadata construction must
accompany the result: variable descriptions are generated with access to
the true graph, although prompts request descriptions that do not
explicitly state the links. The same elicited constraint set is reused
across data repetitions for a dataset. Such repetitions assess sampling
variation conditional on that set; they are not independent replications
of the full elicitation process. Neutral descriptions, alternative
descriptions and misleading semantic controls remain necessary in a
banking prototype. \citep{exec14_expert_argumentation}

Other qualifications affect the comparison. The consensus procedure can
return no constraints and sacrifices recall for precision. Runtime
tables omit external model latency. Hidden-confounder extensions and
calibrated adaptive weights are proposed future work. The appendix's
statistical-test description also uses inconsistent observation and
degrees-of-freedom descriptions, so its significance calculations are
not independently reconstructed here. The proposal can adopt the idea of
explicit conflict traces while keeping the reported gains as bounded
literature evidence and leaving runtime adoption subject to a separate
code, license and environment review.
\citep{exec14_expert_argumentation}

\hypertarget{evidence-adequacy-differs-from-the-truth-of-an-assumption}{%
\subsection{Evidence adequacy differs from the truth of an
assumption}\label{evidence-adequacy-differs-from-the-truth-of-an-assumption}}

ARGUS audits whether a paper reports evidence supporting
difference-in-differences assumptions. It expressly does not adjudicate
causal truth. Its bounded retrieval and assessment stages can locate
design flaws that a keyword baseline misses, but retrieval gating trades
some omission detection for abstention. In the five-paper pilot, it
abstains on 22 of 55 cells; among the 33 answered cells, 25 judgments
are more severe than the reconciled human labels. The subsequent
calibration is evaluated on the same small pilot and does not establish
performance on unseen papers. \citep{v12_la_2609_30867}

This evidence suggests three separate unknown states for the banking
workflow. First, the source may not be accessible or the retrieval
procedure may fail. Second, an inspected source may simply not discuss
the needed assumption. Third, it may state an assumption without
evidence that establishes its applicability to the current accounts,
dates or intervention. An operational manual can support the claim that
a randomization procedure was prescribed; assignment logs and an
appropriate design are still needed to establish which procedure
actually occurred. A paper's favorable finding in another population
does not settle transport to the proposed collection policy.

A claim record should identify the source version and passage, the
population, time and intervention it concerns, what the passage
supports, the contrary evidence and the remaining inference. It should
distinguish direct observation, reported procedure, expert assertion,
mathematical implication and hypothesis. Repeated copies should point to
a common source family. An expert should be able to correct the
extracted claim without silently replacing the original text. The system
should retain both interpretations when the disagreement is substantive,
then show which graph restrictions and estimands each would permit.
These are proposed product requirements; this synthesis does not claim
that an agent already performs them reliably.

\hypertarget{a-banking-contrast-that-documents-alone-cannot-identify}{%
\subsection{A banking contrast that documents alone cannot
identify}\label{a-banking-contrast-that-documents-alone-cannot-identify}}

Let \(A\) indicate receipt of a contact offer at decision time zero, and
let \(Y\) indicate a repayment by day thirty. The eligible population is
fixed before assignment. Suppose the value of repayment is USD 100 and
each offer costs USD 20. The target is the incremental expected value
per eligible account, \(100\{E[Y(1)]-E[Y(0)]\}-20\). Both the
intervention and the cost are stipulated for this illustrative
calculation; there is no observed bank experiment.

Consider two mechanisms. In the first, \(A\) is an independent fair coin
and \(Y=A\mathbin{\oplus}E\), where \(E\) is an independent Bernoulli
variable with probability \(1/4\) and \(\oplus\) denotes exclusive-or.
In the second, a latent pre-offer state \(U\) is a fair coin,
\(A=U\mathbin{\oplus}E\), and \(Y=U\). Both mechanisms produce exactly
the same observed distribution:

\[
\begin{array}{c|cc}
 &Y=0&Y=1\\\hline
A=0&3/8&1/8\\
A=1&1/8&3/8
\end{array}.
\]

They therefore have the same observed repayment contrast of \(1/2\).
Under an intervention, however, the first mechanism gives repayment
probabilities \(1/4\) and \(3/4\), whereas the second gives \(1/2\)
under either action. The causal risk differences are \(1/2\) and zero,
and the net values are USD 30 and minus USD 20. In the restricted class
containing just these two mechanisms, the effect set is \(\{0,1/2\}\)
and the net-value set is \(\{-20,30\}\). These sets are not advertised
as sharp bounds over every possible structural model.
\citep{exec14_expert_oracle}

An assistant that selects the first diagram because several sources call
the offer effective would authorize a different decision from one
retaining the latent-selection explanation. More rows from the same
observed law cannot distinguish these two mechanisms. A useful expert
question concerns the actual assignment rule and the information
available at assignment, rather than a generic opinion about whether
contact helps. A verified randomized assignment or defensible additional
measurement and assumptions can change the inference. Until then,
retaining both models and refusing a uniquely identified benefit is an
informative result.

This example also supplies an inexpensive formal baseline for a
prototype. Enumerate both candidate mechanisms, verify their
observational equality and compare the assistant's admitted graph and
effect claim with the declared information. The correct behavior under
ambiguous evidence is to retain the distinction between observational
fit and intervention value. The exercise evaluates consistency with
stipulated truth, not whether a public document accurately describes an
actual bank.

\hypertarget{why-repeated-agreement-can-exaggerate-evidence}{%
\subsection{Why repeated agreement can exaggerate
evidence}\label{why-repeated-agreement-can-exaggerate-evidence}}

Suppose an illustrative hypothesis has prior probability \(1/2\) and one
independent source has a known likelihood ratio of four. The posterior
is \(4/5\). Five copies of the same source still supply that same
evidence. If they are incorrectly multiplied as independent likelihood
ratios, the reported posterior becomes \(4^5/(1+4^5)=1024/1025\). In
practice a citation rarely comes with a defensible numerical likelihood
ratio at all. The example demonstrates the need to track shared
provenance rather than prescribing a Bayesian score for every passage.
\citep{exec14_expert_oracle}

The analogous issue arises with repeated model calls. For a binary
judgment with independent error probability \(1/5\), five-call unanimity
has probability \((4/5)^5+(1/5)^5=41/125\); the probability of an
incorrect answer conditional on unanimity is \(1/1025\). Now impose
perfectly shared errors with the same marginal error probability. Every
call agrees, and the conditional error remains \(1/5\). These exact
constructed laws show that marginal reliability and agreement alone do
not determine the value of consensus. They do not estimate the
dependence or accuracy of any current model.
\citep{exec14_expert_oracle}

The prototype should accordingly compare individual judgment, consensus,
source-family aggregation and a human-led baseline. It should vary
misleading shared evidence, missing evidence and contradictory passages,
while retaining all requests in the denominator. Conditional correctness
among returned claims must appear beside availability, false assertions
of no relationship and incorrect refusals. Source-family or task-family
dependence also matters for uncertainty: many paraphrases of one
mechanism do not provide the same evidence as many independent
mechanisms.

\hypertarget{what-a-package-can-contribute-and-what-must-be-repaired}{%
\subsection{What a package can contribute, and what must be
repaired}\label{what-a-package-can-contribute-and-what-must-be-repaired}}

The inspected PyWhyLLM revision of 6 January 2026 provides concrete
suggestion and retrieval components under a root MIT notice. Its
packaging uses a dynamic version placeholder and a substantial model,
embedding and retrieval dependency stack; that declaration is not a
qualified environment. The selected identification wrapper delegates
backdoor suggestions to a model-based confounder suggester. Its
frontdoor method is a placeholder, while a different simple class does
generate frontdoor names through a prompt. Neither a method name nor a
generated list establishes that a graphical identification theorem has
been applied. \citep{exec14_expert_pywhy}

The role-based confounder path initializes a treatment-to-outcome edge,
asks model personas for candidate confounders and accumulates counts.
Such counts represent agreement across prompts, with no calibrated
probability or source independence attached. Some selected calls retry a
missing-key error without a finite bound in the inspected loop. These
findings support explicit caller budgets and preservation of failure;
they are static-source findings, not observed failures of an executed
service. \citep{exec14_expert_pywhy}

The augmented pairwise path adds a second boundary. It downloads a
knowledge resource at construction, finds a matching pair through
lexical and embedding retrieval, and concatenates selected source
sentences into the retrieval context. That concatenation does not carry
their document identity into a structured claim record. When no matching
pair is found, it falls back to model knowledge. The answer parser
accepts only the two directions or neither; neither is not an explicit
insufficient-evidence state. The helper also creates its own model
client instead of consistently using the wrapper's supplied model. These
behaviors make direct adoption inappropriate for the proposed evidence
contract. \citep{exec14_expert_pywhy}

A permissible first integration would be a small, optional elicitation
component whose output remains a proposal for review. It needs explicit
model configuration, pre-acquired versioned source inputs, retained
passage identities, bounded calls and retries, and distinct supported,
contradicted, unknown and unavailable outcomes. Graph conversion must
keep negative and unresolved judgments as records, validate clocks and
direction restrictions, and run a separately specified identification
procedure. Dataset licenses, transitive dependencies, provider behavior
and the remaining import paths require review before the component is
enabled. The current review enables no adapter.

The alternatives remain substantive. A manual evidence table with an
explicit diagram may be sufficient for a small proposal. A deterministic
retrieval and graph-checking system can expose evidence without
delegating judgment. Soft priors preserve alternative graphs but require
sensitivity to their strength. Argumentation provides explicit conflict
traces but requires a policy for hard and defeasible restrictions. An
unqualified hard-constraint pipeline can produce confident answers
efficiently by excluding the correct model.

The known next step is therefore a public/generated-data prototype
comparing these workflows under declared evidence conditions. It should
measure total retrieval and model cost, expert correction time, failed
requests and the effect of evidence handling on the banking decision.
The unresolved research concerns calibration across new domains,
dependence among sources and models, and useful uncertainty when graph
and observation assumptions remain disputed. The truth of an
undocumented historical assignment cannot be manufactured by a better
prompt. That boundary is part of the proposed product's behavior, while
bank experiments and institutional acceptance remain later work.
The next study supplies the joint-graph calculation behind this advice: a
fallible expert can narrow a legal graph class only under an explicit model
of errors and a clearly stated retention criterion.

\label{exec14:expert-evidence-end}


% SOURCE docs/release/2026-10-29/source/docs/writing-repair/2026-10-29/src/3c5c615100571a83.tex CHAR 0
\section{Imperfect advice, joint graph evidence and retention guarantees}
\label{exec17:advice-development}

A banking analyst may know that a diagram remains ambiguous while still
having useful information about the disputed arrows. An account officer can
describe which event usually precedes another; a procedure can explain a
permitted action; a language model can suggest which interpretation to
investigate. The difficult question is how much graph uncertainty this advice
can remove without concealing its mistakes. Requiring an adviser to be
infallible discards potentially useful information. Treating each confident
answer as a causal fact discards potentially correct graphs. A proposal needs
an explicit rule between these extremes, together with a statement of the
probability that the rule actually controls.

Long and colleagues formulate this problem using language models as imperfect
experts. Their 2023 first version starts from the true Markov equivalence class
and uses metadata-based judgments to orient its ambiguous edges. A prior over
legal graphs and an adviser-error likelihood produce joint posterior
probabilities; greedy choices trade a smaller class against the modeled risk
of excluding its true member. This is a useful methodological alternative to
hard acceptance of an elicited diagram. The known starting class and the
specified adviser model are substantive premises, however, rather than
consequences of a natural-language interface.~\citep{exec17_imperfect_experts}

\subsection{The starting class and the experiment that supplies advice}

Let $G$ be a directed acyclic graph, and let $\mathcal M_0$ be a finite class of
graphs with the same skeleton and unshielded colliders. For an ambiguous pair
$j$, write $O_j(G)\in\{0,1\}$ for its orientation in $G$ and $E_j$ for the
adviser's proposed orientation. These binary judgments concern an existing
adjacency; they do not decide whether the two variables are the right
measurements, whether an edge is absent, or whether a hidden common cause
should replace it. In particular, the method's pair question does not provide
a general test of the bank's variable ontology or latent-variable assumptions.

Under a uniform prior on $\mathcal M_0$ and a conditionally independent adviser
with constant error $\epsilon$, the likelihood and posterior are
\begin{equation}
\begin{split}
 L(g;e)&=\prod_{j=1}^{u}(1-\epsilon)^{\mathbf 1\{O_j(g)=e_j\}}
                   \epsilon^{\mathbf 1\{O_j(g)\ne e_j\}},\\
 \pi(g\mid e)&=\frac{\mathbf 1\{g\in\mathcal M_0\}L(g;e)}
                        {\sum_{h\in\mathcal M_0}L(h;e)}.
\end{split}
\label{exec17:advice-posterior}
\end{equation}
Here $u$ is the number of ambiguous pairs, $e$ is the observed advice vector,
and $\mathbf 1$ denotes an indicator. The likelihood factorization is a claim
about how adviser errors arise conditional on the graph. The posterior does
not factorize over edges, because the legal graph class couples their
orientations. Long et al. make this distinction explicitly. A cycle or a new
unshielded collider can make a collection of individually plausible answers
jointly impossible.~\citep{exec17_imperfect_experts}

A retained set $\mathcal R(e)\subseteq\mathcal M_0$ can be required to satisfy
\begin{equation}
 \sum_{g\in\mathcal R(e)}\pi(g\mid e)\geq 1-\eta,
 \qquad 0\leq\eta\leq1.
\label{exec17:advice-retention}
\end{equation}
The tolerance $\eta$ describes excluded posterior mass under this model.
The paper considers greedy choices based on the resulting class size or
retention risk, propagating consequences of each chosen orientation. It does
not establish that its greedy rule solves every global minimum-cardinality
retention problem. More fundamentally, the word ``probability'' in
\eqref{exec17:advice-retention} needs its probability model. A posterior over
graphs is useful conditional uncertainty; it is not automatically a
confidence statement for every fixed bank graph and every data-dependent
starting class.

\subsection{A small legal class exposes dependence between answers}

Consider the chain skeleton $A-B-C$, with no unshielded collider. Its three
legal DAGs are $g_1:A\to B\to C$, $g_2:A\leftarrow B\to C$, and
$g_3:A\leftarrow B\leftarrow C$. Code a forward edge as one, so their
orientation vectors are $(1,1),(0,1),(0,0)$. The fourth vector $(1,0)$ would
create $A\to B\leftarrow C$ and is absent from this equivalence class.
For $\epsilon=1/10$, advice $(1,1)$ gives
\begin{equation}
 \bigl(\pi(g_1\mid e),\pi(g_2\mid e),\pi(g_3\mid e)\bigr)
       =\frac{1}{91}(81,9,1).
\label{exec17:advice-chain}
\end{equation}
This calculation follows by normalizing the likelihoods
$(81,9,1)/100$. It is an exact project example, independently enumerated over
the legal graphs, rather than a reported result of the paper's software.
\citep{exec17_advice_validity_exact}

The posterior mass of $A\to B$ is $81/91$, approximately $0.8901$, while that
of $B\to C$ is $90/91$, approximately $0.9890$. At a posterior retention
threshold of $0.9$, the latter orientation can be retained by keeping
$\{g_1,g_2\}$; keeping only $g_1$ fails the threshold. Two answers generated by
an adviser with marginal error $0.1$ have not produced a graph with posterior
probability $0.9$. The legal-class restriction and joint evidence determine
the result. If instead the advice is the incompatible vector $(1,0)$, the
posterior becomes $(9,1,9)/19$. Each advised orientation then has posterior
mass $9/19$. Disagreement with the graph constraints provides a reason to
reduce confidence, not a license to force the fourth, inadmissible graph into
the class.

For a bank, this example illustrates why one cannot independently approve an
``advice causes activity'' arrow and a ``credit draw causes activity'' arrow
and assume that both fit the diagram already justified by data. The variables
need precise definitions and clocks before that particular class is even
plausible. If documented time ordering rules out one orientation, it is
additional structural information and should change the class explicitly.
It should not be disguised as stronger language-model confidence. Conversely,
if a credible account requires a hidden common cause, the three-DAG class is
the wrong model family; redistributing its posterior cannot repair the missing
graph.

\subsection{Repeated prompts and the missing adviser law}

The constant-error model is attractive because it makes the calculation
transparent. Its applicability to a language model is an empirical question.
Long et al. score the two directional answer tokens and normalize their log
probabilities with a softmax. Their fixed prompt gives deterministic output;
the experiments randomize among causal verbs to introduce variation. This
historical procedure concerns text-davinci-002 and text-davinci-003. A softmax
score is a model preference over the presented alternatives, while the needed
likelihood concerns error relative to a true orientation. Agreement between
the two quantities requires evidence.~\citep{exec17_imperfect_experts}

The distinction is especially visible when an assistant asks the same question
repeatedly. Let the true direction $G$ have equal prior probabilities and let
one latent answer $S$ agree with $G$ with probability $0.9$. Suppose five calls
all repeat $S$. The correct joint model gives
\begin{equation}
 \Pr(G=1\mid E_1=\cdots=E_5=1)=0.9,
 \qquad
 \widetilde\Pr(G=1\mid E_1=\cdots=E_5=1)
        =\frac{9^5}{1+9^5}=\frac{59049}{59050}
\label{exec17:advice-repetition}
\end{equation}
for, respectively, the repeated-signal law and an incorrectly factorized
five-signal law. Enumerating the graph and latent-answer joint distribution
reproduces both calculations. The second number is approximately $0.999983$;
no corresponding new information entered the first experiment.
\citep{exec17_advice_validity_exact}

Changing a verb, sampling a different completion or asking a second model does
not by itself establish conditional independence. Models can share training
material, prompts can share a misconception, and several retrieved passages
can repeat one source. A practical adviser model may therefore need grouped
errors, variable-pair difficulty, abstention, asymmetric mistakes and dependence
on the wording or supplied documents. These are proposed modeling extensions,
not properties already verified for the bank. An alternative is to keep
elicited answers as soft sensitivity inputs and report how conclusions change
as their influence varies. Another is to use the answers only to prioritize
human investigation while leaving the admissible graph class intact.

\subsection{Posterior retention and coverage of a learned class}

Suppose the starting class is instead learned from observations $D$, yielding
$\widehat{\mathcal M}(D)$. If the true graph $g_\star$ is absent from that
class, every refinement that returns a subset also excludes it. Even a
perfectly computed adviser posterior cannot recover an omitted member. A
frequentist composition would need an additional argument. For example, if
\begin{equation}
\begin{split}
 \Pr_{g_\star}\{g_\star\notin\widehat{\mathcal M}(D)\}&\leq\alpha,\\
 \Pr_{g_\star}\{g_\star\notin\mathcal R(D,E)
       \mid g_\star\in\widehat{\mathcal M}(D)\}&\leq\eta,
\end{split}
\label{exec17:advice-coverage-premises}
\end{equation}
then the law of total probability gives
\begin{equation}
 \Pr_{g_\star}\{g_\star\notin\mathcal R(D,E)\}
       \leq\alpha+(1-\alpha)\eta\leq\alpha+\eta.
\label{exec17:advice-coverage-composition}
\end{equation}
The second premise concerns the actual joint data/advice experiment, including
selection of the class. It cannot be replaced by a pointwise posterior-mass
calculation. The displayed bound is a project derivation under the two stated
premises, not a new qualification result for the published method or the
proposed library.

This distinction affects downstream monetary decisions. If a policy is chosen
after selecting a graph, uncertainty about the retained graph and uncertainty
about its causal parameters both matter. For a policy $a$, a proposed robust
decision could examine the smallest identified value across retained graphs
and parameter laws. Such a value set must retain any nonidentification within
a graph; a posterior orientation does not supply missing intervention data.
A posterior expected-value decision is a different alternative, conditional
on the chosen graph and parameter priors. The proposal should present those
two decision principles with their assumptions and compare the actions they
recommend, rather than silently turning either into a bank guarantee.

\subsection{What to reuse and what a prototype must still establish}

The published experiments provide a useful warning as well as a useful method.
They derive the starting MEC from each benchmark's true graph, repeat the
synthetic-error experiments five times and the language-model experiments
eight times, and report failures associated with ALARM metadata ambiguity and
text-davinci-003 overconfidence. These findings motivate sensitivity to the
adviser law. They do not evaluate discovery errors in a learned starting class
or establish performance for a current banking ontology and model. The
paper's code and codebook links identify candidates for further inspection;
the paper alone does not establish their license, current dependencies or
runtime behavior.~\citep{exec17_imperfect_experts}

The proposal can reuse the conceptual separation between a data-supported
class and fallible directional advice, together with joint consistency and
retention calculations. It should reject a default that hard-deletes every
graph contradicted by an uncalibrated answer. Direct code adoption or a port
needs a pinned source and rights review, agreement with independent small-class
references, and a declared regime in which enumeration or an approximation is
feasible. A scalable approximate posterior also needs an error statement;
having a graph-search API does not settle its approximation quality.

A public prototype can test these issues without expensive bank experiments.
It can start from stipulated graph classes, preserve independent graph truth,
separate correct-class and misspecified-class cases, vary adviser dependence,
and compare hard constraints, soft priors, abstention and unchanged-class
baselines. Renamed variables and genuinely unfamiliar mechanisms are useful
additional cases, but they test different generalizations. Report graph
retention, class size and downstream action loss separately. Such experiments
would establish whether a particular implementation deserves a banking trial;
the exact examples above establish why those comparisons are necessary.

\label{exec17:advice-development-end}


% SOURCE docs/release/2026-10-29/source/docs/writing-repair/2026-10-29/src/c364c3a2f1bede04.tex CHAR 0
\section{What an elicited Bayesian prior actually buys}
\label{exec18:bayesian-advice}
A bank considering an LLM-assisted deposit diagram needs more than an answer to whether the proposed arrows look plausible. It needs to know which uncertainty distribution the answers define, how that distribution interacts with transaction data, and which observations could discredit the combination. Videira's corrected master's thesis provides a useful concrete case: it combines elicited edge probabilities with DiBS, then considers feedback to a score-based ancestral-graph generator. These are different graph domains and different uses of advice. The thesis was defended on 23 February 2026 and the university records publication on 19 March 2026, before the declared literature cutoff. The comparison here uses the genuine corrected thesis, including its experimental tables and appendices; it does not independently reproduce either upstream system~\citep{exec18_videira_thesis}.

\subsection{Response frequency is an elicitation quantity}
The first construction queries each variable pair ten times and estimates frequencies of three responses: one direction, the other direction, or no relation. The text replaces probabilities one and zero by 0.99 and $10^{-10}$, respectively, then uses products of edge and nonedge factors as a graph prior. It reports 30 starting particles and 3000 DiBS iterations. These choices specify a finite elicitation and fitting procedure; they do not establish that a frequency of 0.9 is a 90\% probability that the causal assertion is true~\citep{exec18_videira_thesis}.

To see the distinction, fix a prompt and a response mechanism which selects one arrow with probability $9/10$. The mechanism can be identical in a world where that arrow is true and a world where it is false. Repeating the prompt estimates the response mechanism in both worlds, without distinguishing their truth. More draws reduce sampling variation in the response-frequency estimate under an appropriate sampling law. They do not identify adviser accuracy, and correlated draws can reduce even that benefit. In a banking proposal, frequency should therefore be stored as an elicitation statistic together with its prompt, model identity and sampling scheme. A calibrated truth channel needs additional evidence on suitably independent labeled questions~\citep{exec18_bayesian_reference}.

The semantic categories also require repair before direct adoption. The thesis's first prompt offers no statistical association as one response and direct causal directions as the alternatives. No direct edge and no statistical association are different statements. For example, an offer and a subsequent balance can be associated through a pre-existing customer-risk variable even if neither directly causes the other. A proposed adapter must define whether it queries adjacency, ancestry, conditional independence, or an intervention effect, then permit outcomes appropriate to that question. A forced three-way answer cannot silently supply all four kinds of information. This is a limitation of the inspected wording, not a newly observed performance result~\citep{exec18_videira_thesis}.

\subsection{Graph admission changes the distribution}
An edge-factor expression is not yet a distribution over admissible causal graphs unless its support and normalizer are specified. For a deliberately small example, let the two possible arrows between two variables be independent Bernoulli indicators with probability $1/2$ each. The unconstrained model gives four configurations equal probability. Rejecting the two-cycle leaves the empty graph and the two single-arrow graphs, each with conditional probability $1/3$. Consequently
\begin{equation}
 \Pr(A\to Y\mid\text{acyclic})=\frac{1/4}{3/4}=\frac13,
 \label{exec18:dag-conditioning}
\end{equation}
although the original edge parameter was $1/2$. The exact reference enumerates these four configurations. A consumer must distinguish the elicited edge parameter from a marginal under the admitted joint graph law~\citep{exec18_bayesian_reference}.

The thesis samples 30 graph realizations for its reported graph metrics and rejects cyclic realizations. That makes the admitted count and the conditional sampling rule important. The inspected account does not establish a distribution-free uncertainty guarantee from those 30 draws. Nor can one assume that independent priors on the two opposing directions preserve a mutually exclusive pairwise response model. A reuse design should record the pair distribution, the graph construction, the legal support, rejection counts and any subsequent renormalization explicitly. The earlier joint-graph discussion explains why individually plausible edges need not form a coherent accepted graph~\citep{exec18_videira_thesis}.

\subsection{The reported comparison is mixed}
The thesis evaluates Sachs, MAGIC-7, MAGIC-ORIGINAL and an authored Summer construction. MAGIC-7 removes genetic variables from the larger dataset and is deliberately treated as a case with omitted confounding; the selected DAG learner's assumptions are therefore an essential part of interpreting its performance. The following are reported posterior mean expected structural Hamming distances, abbreviated E-SHD, from Table 3.2. Lower is better for this metric. They are source-reported means, not new runs or a ranking of general banking suitability~\citep{exec18_videira_thesis}.

\begin{center}
\begin{tabular}{lrrr}
\toprule
Dataset & DiBS baseline & With GPT prior & With DeepSeek prior\\
\midrule
Sachs &23.40&21.70&26.70\\
MAGIC-7 &13.76&16.07&16.24\\
MAGIC-ORIGINAL &153.00&160.05&Not reported\\
Summer &12.17&12.42&9.68\\
\bottomrule
\end{tabular}
\end{center}

The GPT-prior differences relative to baseline are $-1.70$, $2.31$, $7.05$ and $0.25$, respectively. The DeepSeek-prior differences are $3.30$, $2.48$ and $-2.49$ where reported. These exact differences of printed means reproduce separately; they are not significance tests. The printed error bars are not converted here into standard errors or confidence intervals because that interpretation requires their actual experimental denominator and dependence structure~\citep{exec18_bayesian_reference}. On MAGIC-ORIGINAL, the reported mean AUROC also declines from 0.62 to 0.60 with the GPT prior. On Summer, the DeepSeek prior improves both reported metrics. The defensible conclusion is a conditional opportunity with adverse examples, rather than uniform improvement from adding an LLM~\citep{exec18_videira_thesis}.

Two source statements deserve particular care. The discussion calls AUROC 0.5 the worst possible value. A perfectly reversed binary ranking instead has AUROC 0, as the two-observation exact reference demonstrates. The final discussion also describes the combined method as better than the original prior in terms that are broader than the table supports. The proposal follows the tabulated, dataset-specific evidence. It retains the author's concern about known datasets in model training, sparse graphs favoring some SHD comparisons, and violated data-generating assumptions. A favorable prior cannot by itself repair a misspecified likelihood~\citep{exec18_videira_thesis,exec18_bayesian_reference}.

\subsection{Feedback to an ancestral-graph generator has another contract}
The second construction uses an ancestral-graph GFlowNet with a score-based reward, graph-validity constraints and expert feedback. The thesis describes four relationship categories, a confidence parameter for each queried pair, and a categorical channel that places probability $\pi$ on the reported true category and $(1-\pi)/3$ on each alternative. It uses repeated model responses to supply the confidence. Its distributional assessment uses 20 generated five-node graphs with 500 observations each. The separate feedback experiment uses reduced Sachs graphs and reports three simulations in the figure caption~\citep{exec18_videira_thesis}.

For a transparent conditional calculation, take prior category probabilities $(0.4,0.3,0.2,0.1)$, channel accuracy $\pi=0.7$, and feedback selecting category 2. The unnormalized posterior weights are $(0.04,0.21,0.02,0.01)$, summing to 0.28, so
\begin{equation}
 \rho(\cdot\mid f=2)=\left(\frac17,\frac34,\frac1{14},\frac1{28}\right).
 \label{exec18:feedback-update}
\end{equation}
This is a valid update under the stipulated channel. If 0.7 merely measures repeated-answer consistency, the assumed likelihood is not thereby verified. Moreover, pairwise relationship marginals do not determine every dependence in the joint graph distribution. An adapter needs to specify the graph-level update and validate its scientific interpretation~\citep{exec18_bayesian_reference}.

The thesis reports improved E-SHD as feedback accumulates in its selected Sachs experiments, while noting worse unreported DeepSeek trials. Those observations do not prove monotone improvement or convergence to the true graph for arbitrary advisers. The selected reward is a transformation of a fit score; accurate sampling from that reward establishes fidelity to the chosen target, not causal truth or frequentist graph coverage. The relation vocabulary also needs an explicit reconciliation: the mathematical description uses ancestry while parts of the elicitation wording request direct relationships. The source records both descriptions; the proposal does not silently equate them~\citep{exec18_videira_thesis}.

\subsection{A bounded reuse and evaluation decision}
The thesis is valuable as a design and adverse-result source. Its Appendix A supplies an R data-generation listing and Appendix B gives DSPy field declarations. These are not a complete pinned package, trained generator, reproducible prompt transcript or dependency manifest. In particular, the R listing draws random structural parameters before its displayed seed call, and requests 2000 observations while the body describes 2001. Without the missing initialization or saved parameters, that listing alone does not reproduce the exact structural model. A preserving reproduction should retain these discrepancies, specify a seed before every relevant draw, save the realized parameters, and declare its actual sample count as a new experiment~\citep{exec18_videira_thesis}.

The proposed disposition is therefore to adopt the separation of elicitation, graph admission and updating as a design pattern; retain DiBS and ancestral-graph methods as distinct backend candidates; and defer executable integration until exact source versions, code rights, prompts, graph encodings and experiment artifacts are inspected. The thesis's CC BY 4.0 notice supports attribution to its text but is not automatically a license for every external dependency or model service. No downloaded code is executed here.

A useful inexpensive prototype compares a data-only learner, elicited advice alone and the combined method under the same graph law and budget. It should include correct and wrong advice, correlated repeated answers, causal-sufficiency violations, unfamiliar variable names, rejected graph samples and full attempted-query costs. Graph scores should be supplemented with the resulting intervention-query decisions and abstentions. In a deposit application, choosing the wrong adjustment relation can change the monetary target even when only one edge is wrong. The evaluation must preserve that consequence. Authentic bank labels and costly experiments follow only after a scoped public/generated comparison shows merit; no current evidence warrants a universal improvement, calibrated semantic truth or bank-adoption claim.
\subsection{Definitions and advice constrain different parts of the model}

A variable definition establishes which quantity a record is intended to
measure; a measurement model explains its relation to a scientific construct;
an adviser model describes how judgments bear on graph alternatives. None can
silently substitute for another. The examples show why accurate labels can
misstate an effect, repeated evidence can exaggerate confidence, and graph
admission can change the distribution implied by individual edge parameters.

The practical output is a documented set of definitions and restrictions,
with unresolved alternatives and the information needed to distinguish them.
Deterministic dictionaries, human evidence review, soft-prior sensitivity and
formal graph calculations remain substantive alternatives to additional model
calls. The next chapter develops the diagrams and disagreement review that
consume these definitions, preserving their scientific and measurement limits.
\label{exec18:bayesian-advice-end}

